\chapter*{VI/32 Legacy-Problem and Scope Audit}
\addcontentsline{toc}{chapter}{VI/32 Legacy-Problem and Scope Audit}
\label{audit:vi32}

\textbf{Audit marker:} \texttt{VI32\_LOCAL\_GEOMETRY\_INSTANCE\_AUDIT\_2026-08-23}.

\section*{AG8 tangent/local-dimension stream}
The migration atlas routes the AG8 tangent/local-dimension descendants uniquely to VI/32.
The chapter reconstructs that stream through local-ring dimension, $\mathfrak m/\mathfrak m^2$,
Zariski tangent spaces, regularity, singularity, and elementary Jacobian tests.

\section*{AG8 Problem 2 split completed}
Earlier chapters assigned function-field material to VI/27, global/open dimension to VI/30, and codimension to VI/31.
The remaining closed-point local-dimension theorem is now proved as VI.32.P1:
\[
 x\text{ closed in a variety }X\quad\Longrightarrow\quad\dim\OO_{X,x}=\dim X.
\]
This completes the dimension-side routing of the multi-part Problem 2 stream.

\section*{AG8 Problem 3 split completed}
VI/30 consumed the open-dimension failure (part (d)), and VI/31 consumed part (c).
VI/32 consumes the remaining closed-point local-dimension failure, part (a), as VI.32.P2.
The disjoint-union example
\[
 X=\Spec k\sqcup\mathbb A^1_k
\]
has a closed point $x$ on the zero-dimensional component with $\dim\OO_{X,x}=0<1=\dim X$.
Thus no part of AG8 Problem 3 remains reserved after VI/32.

\section*{Legacy DG4 Problem 13 recovered}
The earlier Volume VI audits reserve legacy DG4 Problem 13, the unions-of-lines tangent-space problem,
for VI/32.  It is now preserved as VI.32.P20,
using the two running line-union models already established in VI/05:
\[
 V(xy,xz,yz)\subset\mathbb A^3_k,
 \qquad
 V(uv(u-v))\subset\mathbb A^2_k.
\]
Their component-incidence pictures agree, while their tangent dimensions at the common
three-branch point are $3$ and $2$.  Thus the local tangent invariant separates them.

\section*{Boundary audit}
The chapter uses the intrinsic cotangent space and its dual but does not claim the later theory of Kähler differentials,
smooth morphisms, étale morphisms, completions, or general tangent cones.  It compares regularity with examples of
normalization but does not duplicate the normalization theory of VI/28.
